% Generated draft: requires mathematical and source review.
\documentclass[11pt]{article}
\usepackage[margin=25mm]{geometry}
\usepackage{fontspec}
\setmainfont{Latin Modern Roman}
\usepackage{amsmath,amssymb,mathtools,unicode-math}
\setmathfont{Latin Modern Math}
\usepackage{enumitem,needspace,xurl,hyperref}
\setlength{\parindent}{0pt}
\setlength{\parskip}{5pt}
\setlength{\emergencystretch}{3em}
\newcommand{\N}{\mathbb{N}}
\newcommand{\Z}{\mathbb{Z}}
\newcommand{\Q}{\mathbb{Q}}
\newcommand{\R}{\mathbb{R}}
\newcommand{\C}{\mathbb{C}}
\newcommand{\F}{\mathbb{F}}
\newcommand{\A}{\mathbb{A}}
\newcommand{\PP}{\mathbb{P}}
\newcommand{\E}{\mathbb{E}}
\newcommand{\eps}{\varepsilon}
\newcommand{\dd}{\,\mathrm{d}}
\begin{document}
\section*{UnsolvedMath ID: 700492}
\textbf{Source problem code:} catalogue entry
\subsection{Definitions and mathematical statement}
Let \(G\) be a finite group. A subgroup \(H\leq G\) is soluble if its derived series terminates at the trivial group. Its index is \([G:H]=|G|/|H|\). For a finite soluble group \(H\), define its \(2\)-length \(\ell_2(H)\) to be the minimum number of nontrivial \(2\)-group factors in a normal series
\[
1=H_0\leq H_1\leq\cdots\leq H_r=H,
\]
where every \(H_i\) is normal in \(H\), and every factor \(H_i/H_{i-1}\) is either a \(2\)-group or a group of odd order. Here a \(2\)-group is a finite group whose order is a power of \(2\).

The problem is to classify, up to isomorphism, all finite groups \(G\) such that
\[
\forall H\leq G,\qquad
\bigl(H\text{ is soluble and }[G:H]\text{ is odd}\bigr)
\ \Longrightarrow\ \ell_2(H)=1.
\]
The quantifier includes \(H=G\) whenever \(G\) is soluble. No existence hypothesis for soluble subgroups of odd index is stated; if there are none, the displayed condition holds vacuously.

The equality \(\ell_2(H)=1\) follows the literal source wording. Under the definition above, odd-order groups have \(2\)-length \(0\), so replacing this equality by \(\ell_2(H)\leq1\) would change the problem. The supplied record does not clarify whether such a weaker interpretation was intended.

The source reports that the groups denoted by \(L_n(2^m)\) are known to satisfy the condition. In standard notation these are projective special linear groups \(\operatorname{PSL}_n(\mathbb F_{2^m})\); the record does not explicitly specify parameter ranges or exceptional cases. This example is retained as a source assertion, without independent verification.
\subsection{Short English statement}
Classify the finite groups for which every soluble subgroup of odd index has \(2\)-length exactly \(1\), using the literal wording of the source. The source lists \(L_n(2^m)\) as examples. The exact-versus-at-most interpretation and the example parameters require review.
\subsection{Sources}
\begin{itemize}
\item Local UnsolvedMath registry, permanent ID \(700492\), exact source problem code \texttt{catalogue entry}, title ``Kourovka Notebook Problem 8.42.'' The supplied record labels the problem open; this is a registry classification, not an independently verified assessment. Registry update timestamp: \texttt{2026-10-16T14:00:00Z}.
\item \emph{The Kourovka Notebook (Unsolved Problems in Group Theory)}, original problem 8.42, 8th issue (1982), proposed by A. A. Makhn\"{e}v. The supplied record identifies the 21st-issue edition (2026), arXiv:1401.0300v48, \url{https://arxiv.org/abs/1401.0300v48}. The registry describes the collection as edited by E. I. Khukhro and V. D. Mazurov. The cited edition and the reported family of examples have not been independently checked here.
\end{itemize}
\end{document}
